\chapter{Economic Analysis}\label{chap_EconAnalysis}

\section{Modelling and measuring economic activity}
\subsection{Measures of spatial concentration and specialisation}
The various quantitative measures of economic concentration\index{measures of economic concentration} are defined in detail in this section.

\subsubsection{Locational measures}

\paragraph{Location quotient}
The location quotient\index{location quotient|see{measures of economic concentration}}\index{measures of economic concentration!location quotient} (LQ) is the standard measure of employment distribution that controls for the size of the region. The relative
concentration of industry $i$ in region $j$ is defined as

\begin{eqnarray}
LQ_{i} = \frac{(E_{ij}/E_{in})}{(E_{j}/E_{n})}
\end{eqnarray}

\noindent where, $E_{ij}$ is employment in industry $i$ in region $j$, $E_{j}$ is total employment in region $j$, $E_{in}$ is national employment in industry $i$, and $E_{n}$ is total national employment. Thus, a LQ of greater than one indicates that there is an above average proportion of employment in a given industry in a given region. Industries with an LQ above 1.25 are generally viewed
as constituting the export-oriented economic base.

\paragraph{Horizontal clustering}
Horizontal clustering\index{clustering!horizontal|see{measures of economic concentration}}\index{measures of economic concentration!horizontal clustering} (HC) is an alternative measure that accounts for possible agglomeration effects in terms of the numbers of jobs in a specific industry. As such, it is the arithmetic equivalent to a bubble chart\index{bubble chart} which accounts for an industry's gravitational pull over and above what is expected to occur at the national level. HC can then be written as

\begin{eqnarray}
HC_{i} = E_{ij} - \widehat{E}_{ij},
\end{eqnarray}

\noindent where $\widehat{E}_{ij} = E_{j}\times \frac{E_{in}}{E_{n}}$ by setting $\widehat{LQ}_{i}=1$.

\subsubsection{Shift-share analysis}


\subsubsection{Agglomeration measures}
Location measures generally ignore geographic dispersion and uneven distribution of employment within subregions of an area under study.
Agglomeration measures\index{agglomeration}\index{agglomeration|seealso{measures of economic concentration}} are designed to overcome this shortcoming.

\paragraph{Herfindahl-Hirschman Index}
In the context of industry employment concentration within an MSA the simplest such measure is the Herfindahl-Hirschman index\index{Herfindahl-Hirschman index|see{measures of economic concentration}}\index{measures of economic concentration!Herfindahl-Hirschman index} (HHI). It is defined as

\begin{eqnarray}
HHI_{i} = \sum^{n}_{i=1}s_{i}^{2},
\end{eqnarray}

\noindent where $s_{i}$ is simply industry $i$'s LQ at the level of the subregion compared to the MSA and $n$ is the number of counties
within the MSA.\footnote{In the context of MSA-level data, counties are commonly assumed to be the subregions. For the Blacksburg MSA,
we would set $n=4$ since it includes Giles, Montgomery, Pulaski counties and the city of Radford.} The index is equal to 1 if there is absolute concentration and it takes a value of $\frac{1}{n}$ if employment in the industry is equally dispersed across the MSA.

\paragraph{Locational Gini Coefficient}
The locational Gini coefficient\index{Gini coefficient!locational} (LGC) also accounts for agglomeration and concentration within a specific region, but in a most sophisticated way than the HHI. It is defined as

\begin{eqnarray}
LGC_{i} =
\frac{\sum_{i=1}^{n}\sum_{j=1}^{n}|x_{i}-x_{j}|}{2n(n-1)\mu},
\end{eqnarray}

\noindent where $x_{i}$ and $x_{j}$ are the LQs of industries $i$ and $j$ in each of the subregions, $\mu$ is the mean of the LQs in the MSA and $n$ is the number of counties.

\paragraph{Ellison-Glaeser Index}
The Ellison-Glaeser index\index{Ellison-Glaeser index|see{measures of economic concentration}}\index{measures of economic concentration!Ellison-Glaeser index} (EG) measures how industrial concentration patterns differ from a situation where firms in a purely random manner. The index is defined as follows

\begin{equation}
EG_{i} = \frac{\sum_{i=1}^{n}(s_{i}-x_{i})^2 - (1-
\sum^{n}_{i=1}x^{2}_{i})\sum^{m}_{j=1}z^{2}_{j}}{(1-
\sum^{n}_{i=1}x^{2}_{i})(1- \sum^{m}_{j=1}z^{2}_{j})},
\end{equation}

\noindent where $s_{i}$, $x_{i}$ and $n$ have the same definitions as above; $z_{i}$ is the share of the 6-digit NAICS subsector
establishments in sector $j$ and $m$ is the number of 6-digit subsectors.


\section{Measuring inequality}

\section{Spatial concentration and agglomeration}

\section{Economic impact analysis}
\subsection{Input-output modelling}
\subsection{Applied tools}
\subsubsection{RIMS II}
\subsubsection{IMPLAN}
\subsubsection{FedFIT}

\renewcommand\bibname{Further reading}
\begin{subappendices}
\bibliographystyle{econometrica}
\bibliography{Literature/OverLord}
\IfFileExists{Literature/EconAnalysis.bbl}{\chapter{Economic Analysis}\label{chap_EconAnalysis}

\section{Modelling and measuring economic activity}
\subsection{Measures of spatial concentration and specialisation}
The various quantitative measures of economic concentration\index{measures of economic concentration} are defined in detail in this section.

\subsubsection{Locational measures}

\paragraph{Location quotient}
The location quotient\index{location quotient|see{measures of economic concentration}}\index{measures of economic concentration!location quotient} (LQ) is the standard measure of employment distribution that controls for the size of the region. The relative
concentration of industry $i$ in region $j$ is defined as

\begin{eqnarray}
LQ_{i} = \frac{(E_{ij}/E_{in})}{(E_{j}/E_{n})}
\end{eqnarray}

\noindent where, $E_{ij}$ is employment in industry $i$ in region $j$, $E_{j}$ is total employment in region $j$, $E_{in}$ is national employment in industry $i$, and $E_{n}$ is total national employment. Thus, a LQ of greater than one indicates that there is an above average proportion of employment in a given industry in a given region. Industries with an LQ above 1.25 are generally viewed
as constituting the export-oriented economic base.

\paragraph{Horizontal clustering}
Horizontal clustering\index{clustering!horizontal|see{measures of economic concentration}}\index{measures of economic concentration!horizontal clustering} (HC) is an alternative measure that accounts for possible agglomeration effects in terms of the numbers of jobs in a specific industry. As such, it is the arithmetic equivalent to a bubble chart\index{bubble chart} which accounts for an industry's gravitational pull over and above what is expected to occur at the national level. HC can then be written as

\begin{eqnarray}
HC_{i} = E_{ij} - \widehat{E}_{ij},
\end{eqnarray}

\noindent where $\widehat{E}_{ij} = E_{j}\times \frac{E_{in}}{E_{n}}$ by setting $\widehat{LQ}_{i}=1$.

\subsubsection{Shift-share analysis}


\subsubsection{Agglomeration measures}
Location measures generally ignore geographic dispersion and uneven distribution of employment within subregions of an area under study.
Agglomeration measures\index{agglomeration}\index{agglomeration|seealso{measures of economic concentration}} are designed to overcome this shortcoming.

\paragraph{Herfindahl-Hirschman Index}
In the context of industry employment concentration within an MSA the simplest such measure is the Herfindahl-Hirschman index\index{Herfindahl-Hirschman index|see{measures of economic concentration}}\index{measures of economic concentration!Herfindahl-Hirschman index} (HHI). It is defined as

\begin{eqnarray}
HHI_{i} = \sum^{n}_{i=1}s_{i}^{2},
\end{eqnarray}

\noindent where $s_{i}$ is simply industry $i$'s LQ at the level of the subregion compared to the MSA and $n$ is the number of counties
within the MSA.\footnote{In the context of MSA-level data, counties are commonly assumed to be the subregions. For the Blacksburg MSA,
we would set $n=4$ since it includes Giles, Montgomery, Pulaski counties and the city of Radford.} The index is equal to 1 if there is absolute concentration and it takes a value of $\frac{1}{n}$ if employment in the industry is equally dispersed across the MSA.

\paragraph{Locational Gini Coefficient}
The locational Gini coefficient\index{Gini coefficient!locational} (LGC) also accounts for agglomeration and concentration within a specific region, but in a most sophisticated way than the HHI. It is defined as

\begin{eqnarray}
LGC_{i} =
\frac{\sum_{i=1}^{n}\sum_{j=1}^{n}|x_{i}-x_{j}|}{2n(n-1)\mu},
\end{eqnarray}

\noindent where $x_{i}$ and $x_{j}$ are the LQs of industries $i$ and $j$ in each of the subregions, $\mu$ is the mean of the LQs in the MSA and $n$ is the number of counties.

\paragraph{Ellison-Glaeser Index}
The Ellison-Glaeser index\index{Ellison-Glaeser index|see{measures of economic concentration}}\index{measures of economic concentration!Ellison-Glaeser index} (EG) measures how industrial concentration patterns differ from a situation where firms in a purely random manner. The index is defined as follows

\begin{equation}
EG_{i} = \frac{\sum_{i=1}^{n}(s_{i}-x_{i})^2 - (1-
\sum^{n}_{i=1}x^{2}_{i})\sum^{m}_{j=1}z^{2}_{j}}{(1-
\sum^{n}_{i=1}x^{2}_{i})(1- \sum^{m}_{j=1}z^{2}_{j})},
\end{equation}

\noindent where $s_{i}$, $x_{i}$ and $n$ have the same definitions as above; $z_{i}$ is the share of the 6-digit NAICS subsector
establishments in sector $j$ and $m$ is the number of 6-digit subsectors.


\section{Measuring inequality}

\section{Spatial concentration and agglomeration}

\section{Economic impact analysis}
\subsection{Input-output modelling}
\subsection{Applied tools}
\subsubsection{RIMS II}
\subsubsection{IMPLAN}
\subsubsection{FedFIT}

\renewcommand\bibname{Further reading}
\begin{subappendices}
\bibliographystyle{econometrica}
\bibliography{Literature/OverLord}
\IfFileExists{Literature/EconAnalysis.bbl}{\chapter{Economic Analysis}\label{chap_EconAnalysis}

\section{Modelling and measuring economic activity}
\subsection{Measures of spatial concentration and specialisation}
The various quantitative measures of economic concentration\index{measures of economic concentration} are defined in detail in this section.

\subsubsection{Locational measures}

\paragraph{Location quotient}
The location quotient\index{location quotient|see{measures of economic concentration}}\index{measures of economic concentration!location quotient} (LQ) is the standard measure of employment distribution that controls for the size of the region. The relative
concentration of industry $i$ in region $j$ is defined as

\begin{eqnarray}
LQ_{i} = \frac{(E_{ij}/E_{in})}{(E_{j}/E_{n})}
\end{eqnarray}

\noindent where, $E_{ij}$ is employment in industry $i$ in region $j$, $E_{j}$ is total employment in region $j$, $E_{in}$ is national employment in industry $i$, and $E_{n}$ is total national employment. Thus, a LQ of greater than one indicates that there is an above average proportion of employment in a given industry in a given region. Industries with an LQ above 1.25 are generally viewed
as constituting the export-oriented economic base.

\paragraph{Horizontal clustering}
Horizontal clustering\index{clustering!horizontal|see{measures of economic concentration}}\index{measures of economic concentration!horizontal clustering} (HC) is an alternative measure that accounts for possible agglomeration effects in terms of the numbers of jobs in a specific industry. As such, it is the arithmetic equivalent to a bubble chart\index{bubble chart} which accounts for an industry's gravitational pull over and above what is expected to occur at the national level. HC can then be written as

\begin{eqnarray}
HC_{i} = E_{ij} - \widehat{E}_{ij},
\end{eqnarray}

\noindent where $\widehat{E}_{ij} = E_{j}\times \frac{E_{in}}{E_{n}}$ by setting $\widehat{LQ}_{i}=1$.

\subsubsection{Shift-share analysis}


\subsubsection{Agglomeration measures}
Location measures generally ignore geographic dispersion and uneven distribution of employment within subregions of an area under study.
Agglomeration measures\index{agglomeration}\index{agglomeration|seealso{measures of economic concentration}} are designed to overcome this shortcoming.

\paragraph{Herfindahl-Hirschman Index}
In the context of industry employment concentration within an MSA the simplest such measure is the Herfindahl-Hirschman index\index{Herfindahl-Hirschman index|see{measures of economic concentration}}\index{measures of economic concentration!Herfindahl-Hirschman index} (HHI). It is defined as

\begin{eqnarray}
HHI_{i} = \sum^{n}_{i=1}s_{i}^{2},
\end{eqnarray}

\noindent where $s_{i}$ is simply industry $i$'s LQ at the level of the subregion compared to the MSA and $n$ is the number of counties
within the MSA.\footnote{In the context of MSA-level data, counties are commonly assumed to be the subregions. For the Blacksburg MSA,
we would set $n=4$ since it includes Giles, Montgomery, Pulaski counties and the city of Radford.} The index is equal to 1 if there is absolute concentration and it takes a value of $\frac{1}{n}$ if employment in the industry is equally dispersed across the MSA.

\paragraph{Locational Gini Coefficient}
The locational Gini coefficient\index{Gini coefficient!locational} (LGC) also accounts for agglomeration and concentration within a specific region, but in a most sophisticated way than the HHI. It is defined as

\begin{eqnarray}
LGC_{i} =
\frac{\sum_{i=1}^{n}\sum_{j=1}^{n}|x_{i}-x_{j}|}{2n(n-1)\mu},
\end{eqnarray}

\noindent where $x_{i}$ and $x_{j}$ are the LQs of industries $i$ and $j$ in each of the subregions, $\mu$ is the mean of the LQs in the MSA and $n$ is the number of counties.

\paragraph{Ellison-Glaeser Index}
The Ellison-Glaeser index\index{Ellison-Glaeser index|see{measures of economic concentration}}\index{measures of economic concentration!Ellison-Glaeser index} (EG) measures how industrial concentration patterns differ from a situation where firms in a purely random manner. The index is defined as follows

\begin{equation}
EG_{i} = \frac{\sum_{i=1}^{n}(s_{i}-x_{i})^2 - (1-
\sum^{n}_{i=1}x^{2}_{i})\sum^{m}_{j=1}z^{2}_{j}}{(1-
\sum^{n}_{i=1}x^{2}_{i})(1- \sum^{m}_{j=1}z^{2}_{j})},
\end{equation}

\noindent where $s_{i}$, $x_{i}$ and $n$ have the same definitions as above; $z_{i}$ is the share of the 6-digit NAICS subsector
establishments in sector $j$ and $m$ is the number of 6-digit subsectors.


\section{Measuring inequality}

\section{Spatial concentration and agglomeration}

\section{Economic impact analysis}
\subsection{Input-output modelling}
\subsection{Applied tools}
\subsubsection{RIMS II}
\subsubsection{IMPLAN}
\subsubsection{FedFIT}

\renewcommand\bibname{Further reading}
\begin{subappendices}
\bibliographystyle{econometrica}
\bibliography{Literature/OverLord}
\IfFileExists{Literature/EconAnalysis.bbl}{\chapter{Economic Analysis}\label{chap_EconAnalysis}

\section{Modelling and measuring economic activity}
\subsection{Measures of spatial concentration and specialisation}
The various quantitative measures of economic concentration\index{measures of economic concentration} are defined in detail in this section.

\subsubsection{Locational measures}

\paragraph{Location quotient}
The location quotient\index{location quotient|see{measures of economic concentration}}\index{measures of economic concentration!location quotient} (LQ) is the standard measure of employment distribution that controls for the size of the region. The relative
concentration of industry $i$ in region $j$ is defined as

\begin{eqnarray}
LQ_{i} = \frac{(E_{ij}/E_{in})}{(E_{j}/E_{n})}
\end{eqnarray}

\noindent where, $E_{ij}$ is employment in industry $i$ in region $j$, $E_{j}$ is total employment in region $j$, $E_{in}$ is national employment in industry $i$, and $E_{n}$ is total national employment. Thus, a LQ of greater than one indicates that there is an above average proportion of employment in a given industry in a given region. Industries with an LQ above 1.25 are generally viewed
as constituting the export-oriented economic base.

\paragraph{Horizontal clustering}
Horizontal clustering\index{clustering!horizontal|see{measures of economic concentration}}\index{measures of economic concentration!horizontal clustering} (HC) is an alternative measure that accounts for possible agglomeration effects in terms of the numbers of jobs in a specific industry. As such, it is the arithmetic equivalent to a bubble chart\index{bubble chart} which accounts for an industry's gravitational pull over and above what is expected to occur at the national level. HC can then be written as

\begin{eqnarray}
HC_{i} = E_{ij} - \widehat{E}_{ij},
\end{eqnarray}

\noindent where $\widehat{E}_{ij} = E_{j}\times \frac{E_{in}}{E_{n}}$ by setting $\widehat{LQ}_{i}=1$.

\subsubsection{Shift-share analysis}


\subsubsection{Agglomeration measures}
Location measures generally ignore geographic dispersion and uneven distribution of employment within subregions of an area under study.
Agglomeration measures\index{agglomeration}\index{agglomeration|seealso{measures of economic concentration}} are designed to overcome this shortcoming.

\paragraph{Herfindahl-Hirschman Index}
In the context of industry employment concentration within an MSA the simplest such measure is the Herfindahl-Hirschman index\index{Herfindahl-Hirschman index|see{measures of economic concentration}}\index{measures of economic concentration!Herfindahl-Hirschman index} (HHI). It is defined as

\begin{eqnarray}
HHI_{i} = \sum^{n}_{i=1}s_{i}^{2},
\end{eqnarray}

\noindent where $s_{i}$ is simply industry $i$'s LQ at the level of the subregion compared to the MSA and $n$ is the number of counties
within the MSA.\footnote{In the context of MSA-level data, counties are commonly assumed to be the subregions. For the Blacksburg MSA,
we would set $n=4$ since it includes Giles, Montgomery, Pulaski counties and the city of Radford.} The index is equal to 1 if there is absolute concentration and it takes a value of $\frac{1}{n}$ if employment in the industry is equally dispersed across the MSA.

\paragraph{Locational Gini Coefficient}
The locational Gini coefficient\index{Gini coefficient!locational} (LGC) also accounts for agglomeration and concentration within a specific region, but in a most sophisticated way than the HHI. It is defined as

\begin{eqnarray}
LGC_{i} =
\frac{\sum_{i=1}^{n}\sum_{j=1}^{n}|x_{i}-x_{j}|}{2n(n-1)\mu},
\end{eqnarray}

\noindent where $x_{i}$ and $x_{j}$ are the LQs of industries $i$ and $j$ in each of the subregions, $\mu$ is the mean of the LQs in the MSA and $n$ is the number of counties.

\paragraph{Ellison-Glaeser Index}
The Ellison-Glaeser index\index{Ellison-Glaeser index|see{measures of economic concentration}}\index{measures of economic concentration!Ellison-Glaeser index} (EG) measures how industrial concentration patterns differ from a situation where firms in a purely random manner. The index is defined as follows

\begin{equation}
EG_{i} = \frac{\sum_{i=1}^{n}(s_{i}-x_{i})^2 - (1-
\sum^{n}_{i=1}x^{2}_{i})\sum^{m}_{j=1}z^{2}_{j}}{(1-
\sum^{n}_{i=1}x^{2}_{i})(1- \sum^{m}_{j=1}z^{2}_{j})},
\end{equation}

\noindent where $s_{i}$, $x_{i}$ and $n$ have the same definitions as above; $z_{i}$ is the share of the 6-digit NAICS subsector
establishments in sector $j$ and $m$ is the number of 6-digit subsectors.


\section{Measuring inequality}

\section{Spatial concentration and agglomeration}

\section{Economic impact analysis}
\subsection{Input-output modelling}
\subsection{Applied tools}
\subsubsection{RIMS II}
\subsubsection{IMPLAN}
\subsubsection{FedFIT}

\renewcommand\bibname{Further reading}
\begin{subappendices}
\bibliographystyle{econometrica}
\bibliography{Literature/OverLord}
\IfFileExists{Literature/EconAnalysis.bbl}{\input{Literature/EconAnalysis.bbl}}{\typeout{Need to run bibtex}}
\end{subappendices} }{\typeout{Need to run bibtex}}
\end{subappendices} }{\typeout{Need to run bibtex}}
\end{subappendices} }{\typeout{Need to run bibtex}}
\end{subappendices} 